\section{Results}
\label{sec:results}

We report the completed Codex runs in the inventory generated on 24 September 2026,
separating the broader GPT-5.6 Luna evaluation from the recent ablations. The inventory
counts completed attempts, including unsuccessful submissions and stopping-limit
terminations. For the ablations, all 360 planned sample--configuration pairs enter
the denominator. Execution errors and omissions after blocking errors count as
unsuccessful evaluations, with causes marked separately. We retain the first completed
attempt per sample and configuration when available; retries do not add denominator entries.
% Source: workspace/RESULTS/run-inventory/report.md (2026-09-24T23:13:33.275Z).
% Generated with workspace/report_run_inventory.mjs; per-attempt provenance is in report.json.
% The prior Results section and its editorial notes are preserved in
% sections/09-results-previous-snapshot.tex, but are not included in the paper.

\subsection{Broader evaluation}
\label{sec:results-headline}

In the earlier GPT-5.6 Luna runs at low effort, solve rates rise with assistance, from 4/67
attempts with basic tools to 31/272 with legal enumeration and 28/108 with on-demand target
comparison (Table~\ref{tab:headline}, Appendix~\ref{app:supplementary-results}). These
conditions differ in coverage, repetitions, and budgets, so we treat the trend as descriptive.

\subsection{Tool and reasoning-effort ablations}
\label{sec:results-tiers}
\label{sec:results-ablations}

The recent study targets ten samples from each of four groups per configuration:
easy, medium, hard, and some-layers. It compares GPT-5.6 Luna at low effort with
GPT-6 Luna at low and high effort. The three tool conditions are basic tools,
legal-fold enumeration, and enumeration with automatic tier-3 target feedback.
The run settings are specified in Section~\ref{sec:exp-protocol}. Of 360 planned sample--configuration pairs, 355 have completed episodes,
three have observed execution errors without a completed replacement, and two were
omitted from the restarted batch after blocking errors. All five count as unsuccessful
evaluations under the planned protocol (Appendix~\ref{app:protocol-failures}).

\begin{table}[htbp]
\centering
\caption{Recent Luna ablations. Cells show solved/evaluated samples; each group
has ten intended samples. B: basic tools; L: legal enumeration; A: legal enumeration
with automatic tier-3 target feedback. $*$ includes one input-size error;
$\dagger$ includes one service prompt rejection; $\ddagger$ includes two
omissions after blocking errors. All five count as unsuccessful evaluations.}
\label{tab:luna-ablations}
\small
\setlength{\tabcolsep}{4pt}
\begin{tabular}{lllrrrrr}
\hline
Model & Effort & Tools & Easy & Medium & Hard & Some-layers & Total \\
\hline
GPT-5.6 Luna & Low & B & 2/10 & 0/10 & 0/10 & 0/10 & 2/40 \\
 & Low & L & 1/10 & 0/10 & 0/10 & 5/10 & 6/40 \\
 & Low & A & 4/10 & $0/10^{*}$ & 0/10 & 6/10 & $10/40^{*}$ \\
\hline
GPT-6 Luna & Low & B & 0/10 & 0/10 & 0/10 & 0/10 & 0/40 \\
 & Low & L & 1/10 & $0/10^{*\dagger}$ & 0/10 & 1/10 & $2/40^{*\dagger}$ \\
 & Low & A & 3/10 & $0/10^{\ddagger}$ & 0/10 & 5/10 & $8/40^{\ddagger}$ \\
\hline
GPT-6 Luna & High & B & 3/10 & 0/10 & 0/10 & 3/10 & 6/40 \\
 & High & L & 3/10 & 0/10 & 0/10 & 7/10 & 10/40 \\
 & High & A & 3/10 & 0/10 & 0/10 & 8/10 & 11/40 \\
\hline
\end{tabular}
\end{table}

\paragraph{Reasoning effort.} Pooled over the three tool conditions, GPT-6 Luna solves 6/30
some-layers samples at low effort and 18/30 at high effort. Per condition the increases are
0/10 to 3/10, 1/10 to 7/10 and 5/10 to 8/10; with ten samples per cell only the second is
significant on its own (two-sided Fisher exact test, $p\approx0.02$; $p\approx0.21$ and $0.35$
for the others), and a 95\% Wilson interval spans roughly $\pm0.3$ (e.g.\ 0.11--0.60 for
3/10). We therefore read the effect of effort as consistent across conditions but individually
underpowered. It does not extend to the easy group, where high-effort GPT-6 Luna stays at 3/10
in every condition and low-effort GPT-5.6 Luna with automatic feedback reaches 4/10.

\paragraph{Depth.} No Luna ablation solves a medium or hard sample, and the breadth-first
baseline solves 1/40 medium and 0/16 hard attempts (Appendix~\ref{app:bfs-analysis}), against
366/669 on the easy group. Beyond ten reference folds the benchmark is currently out of reach
for these models and nearly so for systematic search. Depth is not isolated from other geometric differences
between groups.

Supplementary exploratory model results are reported in
Appendix~\ref{app:supplementary-results}.

\subsection{Crease-pattern distance}
\label{sec:results-distance}

A separate scan of 1,421 saved attempts with recoverable states gives a mean crease-pattern
insertion/deletion distance of 45.5, or 0.55 of the target crease count. These averages pool
models, baselines, configurations and repeated samples, and predate the ablations; per-model
values are in Appendix~\ref{app:supplementary-results}.
% EDITING NOTE: Refresh from cp-distance/summary.json after scoring the new runs.

\subsection{Failure analysis}
\label{sec:results-failures}

Rejected actions never end an episode, so failure is a matter of search rather than of
rule-following. In the earlier 447-attempt GPT-5.6 Luna snapshot, 167 episodes submitted a
sequence (63 of them correct), 215 stopped after cycling through previously visited states,
18 stopped on repetition and 47 exhausted the turn budget. Most unsuccessful episodes (215 of
384) therefore end with the model returning to states it has already seen instead of committing
to a new branch.
% EDITING NOTE: Recompute termination counts by cohort and condition from saved traces
% (workspace/analysis/termination_summary.py) and fill the appendix table.

\phantomsection
\label{sec:results-difficulty}
\phantomsection
\label{sec:results-threats}
The broader evaluation, ablations, BFS snapshot and CP-distance scan cover different attempt
sets with different budgets; we report their denominators separately and do not pool them into
a single controlled comparison.
% EDITING NOTE: The old coupling table and main-result plot are preserved in their
% source files and 09-results-previous-snapshot.tex; update their cohorts before reuse.
